\begin{abstract}
% A = Attention getter (motivation)
After major typhoons and heavy rain in Taiwan, average wholesale vegetable
prices have risen by about 10--40\% in the first week and by up to about
100\% in the third week, yet shoppers seeing today's shelf price cannot tell
whether it will rise or fall.
% B = But (challenge)
Market averages are not store prices, retail prices are sticky but jump
abruptly, and large language models (LLMs) asked about prices may produce
unsupported numbers.
% C = Cure (proposed solution)
We propose the Fresh Market Price Forecast Agent (FMPA), a domain AI agent
that anchors short-term forecasts on the store price the user observes.
% D = Development (method)
FMPA restricts the LLM to intent parsing and explanation, obtains every number
through Model Context Protocol (MCP) tools, and estimates the relative price
change at 7 and 14 calendar days with RandomForest and HistGradientBoosting
models, applied to the store price (mode A) and, when a comparable same-store
price from the previous week exists, adjusted by a fixed rule (mode B).
% E = Experiments (evaluation)
We evaluate the forecasting core on one year of daily prices of three
vegetables in 13 public markets in Taichung City with a leakage-safe
chronological split (2,085 seven-day and 1,859 fourteen-day test samples),
against no-change and momentum baselines, using direction macro-F1 and the
mean absolute error of the relative change.
% F = Findings
The models reached macro-F1 of 0.422 and 0.417 but did not outperform the
momentum baseline, and their error exceeded that of the no-change baseline,
consistent with sticky retail prices. Future work will add event indicators
and uncertainty estimates and validate forecasts against matured same-store
prices.
\end{abstract}

\begin{IEEEkeywords}
LLM agent, Model Context Protocol, vegetable price forecasting, sticky prices,
forecasting baselines, decision support
\end{IEEEkeywords}
